% !TeX root = ../main.tex
% Add the above to each chapter to make compiling the PDF easier in some editors.

\chapter{Introduction}\label{chapter:introduction}

\section{Motivation}
Four in every five bugs have been introduced by recent code changes (regeression)~\cite{zhuRegressionGreyboxFuzzing2021b}

Corpus Quality influence on fuzzing
% Focus: Works studying why initial seed quality, corpus distillation, and high initial branch coverage give downstream greybox mutators (like AFL++) a headstart. 
Overall to show: Corpus quality has a significant impact on fuzzing result. (Also explain what good quality means, and name studies that show that these qualities improve fuzzing)

First: Corpus influences fuzzing performance
- EvalFuzzTesting~\cite{kleesEvaluatingFuzzTesting2018} (Study investigated that experimental setups of most studies are flawed and names that a big factor is the initial seed selection, and then do experiments to show so)

Minimal Corpus is Good:
- Corpus Distillation for Effective Fuzzing \cite{herreraCorpusDistillationEffective2020} -> Distill Term!! (Study that compares full, distilled and "empty" (only header present) and comes to conclusion that for big corpora minimizing them is always good)
- Seed Selection for Successful Fuzzing \cite{herreraSeedSelectionSuccessful2021}

Targeted Corpus (Seed selection) is good:
ContinousFuzzingSAP, Bugoss~\cite{yooImprovingConfigurabilityUnitlevel2021}~\cite{kimBugOssBenchmarkRealworld2024} (Select based on functional Coverage)
DART corpus ~\cite{canakciTargetFuzzUsingDARTs2022} (Distance-based selection to target site of seeds)


I want to reframe things a  bit.

The setting is finding bugs introduced by commit changed code because here most bugs are introduces. That is best done before introducing the commit into live code (CI/CD) flow.

All methods are essentially about generating input for a target program until we find a crash.

Mutation Fuzzers do so by taking a set reference input and mutating them until we find a bug. Without starting with such seeds fuzzing is not very effective because we never get deep into the code. They need that "head" start.

They assume that these seeds exisit  new ones based on that.  But when they exisit fuzzing is highly effective at revealing bugs that are hard to find otherwise. More advanded version of fuzzing have then looked into mutationing the seed towards commit changed or later additionaly affected code to trigger regression bugs faster. 



Now that assumption that we have such corpus out of thin air is faulty because the seeds have to be created first by someone or something in the first place. Related work has looked into generation of such seeds out of thin air via various methods and then use the effective mutation fuzzing idea afterwards (either with traditional approaches or using customized solutions).



Our assumption is that we give an even better headstart to mutation fuzzing by creating seeds that already cover as much of affected code as possible by using the semantic capablities of LLMs to understand code (which a fuzzer cant't) and then use a fuzzer afterwards to mutatate (because doing this with LMM would very ineffective).  Showing the creation out of thin air is something we do during evaluation with our tool. One step futher is then to think that we alteratilvy have already a seed corpus generated via one of the various methods, but we want to improve it, but generating additional seeds that are targeted towards these affected areas (because perhaps the exising corpus is missing a seed that covers a certain kind of input format or subformat) or slimming it down first to the most relevenant seeds and then generate addiotonal seeds to generate a more overall targeted corpus (there is work that shows that did this prefiltering of seeds to make them more targeted and showed great sucess with it).